実験結果を書く．結果が出てない場合はやってきたことを書く．
    
\subsection{評価用データセット}
% 以下のデータセットを評価用に利用した．

\subsubsection{データセットの詳細1}
葉っぱの画像を使いました．

\subsubsection{データセットの詳細2}
皮膚がんの画像を使いました．
        
\subsection{実験1}
\subsubsection{実験1-1の詳細}
\input{tables/tex/exp01}

実験1-1の結果を表 \ref{tab:result1}に示す．
表は table generator\footnote{\url{https://www.tablesgenerator.com/}} で作ると非常に楽に作れます。
表の数値は基本的に右揃えなことに注意しましょう。
桁が揃って見やすいです。
\texttt{bookstab} style を使うと最近の論文で使われている表のスタイルになっていい感じになります。
作った表は \texttt{tables/tex/xxxx.tex} のようなファイルを作って貼り付け、\verb|\input{tables/tex/xxxx.tex}| とすることで表示できます。

\subsubsection{実験1-2の詳細}
実験1-2の結果を表 \ref{tab:result2}に示す．

\input{tables/tex/exp02}

\subsection{実験2}
\subsubsection{実験2}
実験2の結果を表 \ref{tab:result3}に示す．

\input{tables/tex/exp03}
    
    